\section{Cubic supports fixed by a two-power diagonal matrix}\label{sec:support}

Let $F\in\C[x_1,\dots,x_n]$ be a cubic form and $B$ the integer matrix whose
columns are the exponent vectors of its monomials with nonzero coefficient
(its \emph{support}). We say $F$ satisfies the \emph{span condition} if
$\one=(1,\dots,1)^t$ lies in the rational column span of $B$.

\begin{lemma}\label{lem:span}
Every cubic form at which some $SL_n$-invariant of positive degree does not
vanish satisfies the span condition, in every coordinate basis.
\end{lemma}

\begin{proof}
Some coefficient monomial of the invariant is nonzero at $F$ and so involves
only coefficients of supported monomials. Invariance under the diagonal torus
of $SL_n$ makes the sum of the corresponding exponent vectors, counted with
multiplicity, a vector with all coordinates equal, and positive since the
degree is positive.
\end{proof}

If a diagonal matrix $A$ with eigenvalues $\alpha_1,\dots,\alpha_n$ of finite
two-power order fixes $F$, the character $\psi:\Z^n\to\C^*$,
$u\mapsto\prod_i\alpha_i^{u_i}$, kills the lattice generated by the support,
and $\psi(\one)=\det A$. If $K$ is a set of support columns forming a basis of
the rational support span, the equation $Kz=\one$ has a unique rational
solution; by Cramer's rule applied to a nonzero maximal minor $K_I$, the least
$d\ge1$ with $dz$ integral divides $\det K_I$, and $d\one=K(dz)$ is in the
support lattice. Hence
\begin{equation}\label{eq:cramer}
 \ord(\det A)\le2^{v_2(\det K_I)} .
\end{equation}

\paragraph{The support graph.}
A monomial $x_i^2x_j$ ($i\ne j$) gives an edge between $i$ and $j$ with
coefficient two at $i$ and one at $j$; a cube $x_i^3$ gives a loop with
coefficient three; all $n$ vertices are included. Every other monomial is of
the form $x_ix_jx_k$ with three distinct indices (a \emph{triple}). For a
spanning tree of a connected component with $v$ vertices, the $v\times(v-1)$
matrix $T$ of its columns has rank $v-1$ (remove a leaf: its row has a single
nonzero entry, one or two). So a graph component has rank $v$ or $v-1$ (an
isolated vertex has rank one if its cube occurs, zero otherwise). Deleting a
row $r$ of $T$ and removing leaves other than $r$ one at a time shows
\begin{equation}\label{eq:treeminor}
 \det T^{(r)}=\pm2^{e_r},\qquad 0\le e_r\le v-1 .
\end{equation}
The signed cofactor vector $u$ of $T$ spans the kernel of $T^t$; along an edge
with coefficients two and one it satisfies $2u_i+u_j=0$, so adjacent entries
differ in absolute value by a factor two and have opposite signs:
\begin{equation}\label{eq:cofactor}
 u=\varepsilon2^{e}\bigl((-2)^{h_v}\bigr)_v,\qquad
 \varepsilon=\pm1,\quad e\ge0,\quad h_v\ge0,\quad\min_vh_v=0 .
\end{equation}

\begin{proposition}\label{prop:graphsupport}
Suppose the support of $F$ consists of cubes and monomials $x_i^2x_j$ only, and
$F$ satisfies the span condition. If a diagonal matrix $A$ of finite two-power
order fixes $F$, then $\ord(\det A)\le2^{n-1}$. If equality holds, then after
a permutation of the variables the support is exactly
\begin{equation}\label{eq:chainsupport}
 x_1^2x_2,\ x_2^2x_3,\ \dots,\ x_{n-1}^2x_n,\ x_n^3,
\end{equation}
the eigenvalues of $A$ are $\xi,\xi^{-2},\xi^4,\dots,\xi^{(-2)^{n-1}}=1$ with
$\xi$ of order $2^{n-1}$, and a diagonal substitution turns $F$ into the chain
cubic $F_n=x_1^2x_2+\dots+x_{n-1}^2x_n+x_n^3$.
\end{proposition}

\begin{proof}
The support lattice, the span condition and $\psi$ split over the connected
components of the support graph, and $\det A$ is the product of the values of
$\psi$ on the all-ones vectors $\one_v$ of the components. Each value has
two-power order; we show it is at most $2^{v-1}$ for a component with $v$
vertices, and at most $2^{v-2}$ if $v\ge2$ and the component is rank-deficient.
An isolated vertex must carry a cube (span condition), and $3e_i$ in the
lattice makes $\psi(e_i)$ trivial.

\emph{Rank $v-1$.} The span of the component is that of a spanning tree $T$,
and $Tc=\one_v$ has a unique rational solution. By Cramer's rule for every
$T^{(r)}$, the denominator $q$ of $c$ divides every $2^{e_r}$, and
$q\one_v=T(qc)$ lies in the lattice. By \eqref{eq:cofactor} adjacent $e_r$
differ by one; all are at most $v-1$, so their minimum is at most $v-2$, and
$q\le2^{v-2}$.

\emph{Rank $v$.} One further column extends $T$ to a nonsingular square
matrix $C$, whose graph is connected and unicyclic (the cycle may be a loop or
a double edge). The order of $\psi(\one_v)$ divides $|\det C|$, so it is at
most $2^{v_2(\det C)}$. Expanding along leaf rows contributes a factor one or
two for each of the $v-l$ tree edges off the cycle of length $l$. A loop
contributes the odd factor three, so $v_2(\det C)\le v-1$. A cycle of length
$l\ge2$ has determinant $\pm(2^a+(-1)^{l-1}2^{l-a})$ for some $0\le a\le l$;
it is nonzero, so $a\ne l-a$ when $l$ is even, and its valuation is
$\min(a,l-a)\le\lfloor(l-1)/2\rfloor$. Then $v_2(\det C)\le v-l+
\lfloor(l-1)/2\rfloor\le v-2$.

The product of two-power orders has order at most the largest one, and every
component has at most $n$ vertices; this gives $\ord(\det A)\le2^{n-1}$. In the
case of equality, some component has $n$ vertices, so the graph is connected,
of full rank, and the chosen unicyclic basis has a loop, at a vertex we call
the root; and every leaf-expansion factor is two. Rooting the tree at the
loop, each tree edge has coefficient two at its child and one at its parent,
and the eigenvalue relations are $\alpha_{\rm root}^3=1$ and
$\alpha_{\rm parent}=\alpha_{\rm child}^{-2}$. As the eigenvalues have two-power
order, $\alpha_{\rm root}=1$, and a vertex at depth $d$ has eigenvalue of order
at most $2^d$; if the depth of the tree is $h$, then $\det A$ has order at most
$2^h$. So $h\ge n-1$, and the tree is a path rooted at one end. The deepest
eigenvalue $\xi$ has order $2^{n-1}$, and the eigenvalues along the path are
$\xi,\xi^{-2},\dots,\xi^{(-2)^{n-1}}=1$. In particular $A$ has order $2^{n-1}$.

These eigenvalues fix no other cubic monomial. Indeed, the exponents
$(-2)^i\bmod2^{n-1}$, $0\le i\le n-2$, have distinct two-adic valuations and
the last exponent is $0$. In a sum of three of them that vanishes modulo
$2^{n-1}$, the least valuation among the nonzero ones cannot occur exactly
once, nor three times (three equal nonzero terms have the valuation of one of
them, since three is odd); so it occurs twice, at one index $i$, and the third
term is $-2(-2)^i$, the next chain exponent, or zero if $i=n-2$. Three zero
terms give $x_n^3$. So the monomials fixed by $A$ are exactly
\eqref{eq:chainsupport}, and they all occur, being the columns of the basis.
Rescale $x_n$ by a cube root and then each $x_i$, $i=n-1,\dots,1$, by a square
root to make all coefficients one.
\end{proof}

\begin{theorem}\label{thm:fullcubicsupport}
Let $F$ be a cubic form in $n\ge2$ variables satisfying the span condition, and
let a diagonal matrix $A$ of finite two-power order fix $F$. Then
$\ord(\det A)\le2^{n-1}$. If equality holds, then after a permutation of the
variables the support of $F$ is exactly \eqref{eq:chainsupport}, $A$ has order
$2^{n-1}$, and a diagonal substitution turns $F$ into $F_n$. The bound is
attained by $F_n$.
\end{theorem}

\begin{proof}
Let $G$ be the matrix of the graph columns (cubes and $x_i^2x_j$). By the
spanning-tree argument, if the graph has $f$ components of full rank and $c$
rank-deficient components, with $v_i$ vertices, then $\rank G=n-c$. In each
full component choose a square basis $C_i$ (a spanning tree and one more graph
column); the proof of Proposition~\ref{prop:graphsupport} gives
\begin{equation}\label{eq:fullcomp}
 v_2(\det C_i)\le v_i-1 .
\end{equation}
In each deficient component choose a spanning tree $T_i$; its signed cofactor
vector is $2^{e_i}u_i$ with $(u_i)_v=\varepsilon_i(-2)^{h_{iv}}$ as in
\eqref{eq:cofactor} (for an isolated vertex without cube, $e_i=h_{iv}=0$).
Put $b_i=v_i-1-e_i$ and $H_i=\max_vh_{iv}$. Since a tree minor is a product of
$v_i-1$ factors one or two, $0\le H_i\le b_i$; and if $H_i=b_i$, the vertex of
height $H_i$ is unique: rooting the tree there, its deleted-row minor has
absolute value $2^{v_i-1}$, so every factor is two, each edge has coefficient
two at its child and one at its parent, $h_{\rm parent}=h_{\rm child}+1$, and
heights decrease strictly away from the root.

Normalize the entries of $u_i$ by $2^{b_i}$. For sums over distinct vertices
$E$ of one deficient component,
\begin{equation}\label{eq:smallsums}
 \Bigl|\sum_{v\in E}\frac{(u_i)_v}{2^{b_i}}\Bigr|\le1,\ \tfrac54,\ \tfrac32
 \qquad\text{for }|E|=1,2,3 .
\end{equation}
If $H_i<b_i$, all entries have absolute value at most $1/2$. If $H_i=b_i$,
there is one entry of absolute value one, which we make positive; the entries
of absolute value $1/2$ are at height $H_i-1$ and so negative; all others have
absolute value at most $1/4$. A pair avoiding the top entry has sum at most one
in absolute value, a pair containing it at most $5/4$; a triple avoiding it at
most $3/2$; a triple containing it gives $0$ if both others are $-1/2$, at
most $3/4$ if one is, and at most $1+1/4+1/4$ otherwise.

The bases $C_i$ and trees $T_i$ form a rational basis of the graph columns.
Extend it by $q$ triple columns $\beta_1,\dots,\beta_q$ to a basis of all
support columns; then $0\le q\le c$ and the support rank is $r=n-c+q$. If
$q=0$, the graph part satisfies the span condition by itself, and
Proposition~\ref{prop:graphsupport} applies: it gives the bound, and in case
of equality the chain eigenvalues, which fix no triple monomial; so $F$ has no
triple monomials and the theorem holds. Assume $q>0$.

Define the $c\times q$ matrix $D_{ij}=u_i^t(\beta_j|_i)$, where $\beta_j|_i$ is
the restriction to the vertices of deficient component $i$. The map sending a
vector to its $c$ pairings with the $u_i$ is surjective with kernel the
rational span of the graph columns (on a full component the graph columns span
all coordinates; on a deficient one $\ker u_i^t$ is the tree span), so $D$ has
rank $q$; choose $q$ deficient components $J$ with $\det D_J\ne0$. Let $K$ be
the $n\times r$ basis matrix. Keep all rows of full components and of the
components in $J$; in every other deficient component delete one row at a
vertex of height zero. The resulting $r\times r$ minor satisfies
\begin{equation}\label{eq:minorfactor}
 \det K_I=\pm\Bigl(\prod_{\text{full }i}\det C_i\Bigr)
 2^{\sum_{\text{deficient }i}e_i}\det D_J .
\end{equation}
Indeed, in a full component eliminate the triple entries with its graph
columns and factor out $\det C_i$; in a deficient component outside $J$ the
retained tree block is square with determinant $\pm2^{e_i}$ (the deleted row
had height zero); eliminate and factor out. These operations stay inside
components. What remains are the tall tree blocks of $J$ and the $q$ triple
columns; expanding along maximal minors of each tree block leaves one row per
component, given by the cofactor vector $2^{e_i}u_i$, and the remaining
$q\times q$ determinant is $\det D_J$.

Divide row $i$ of $D_J$ by $2^{b_i}$. Each column then involves at most three
distinct vertices; by \eqref{eq:smallsums} its squared Euclidean norm is at
most $3$ (three components), $(5/4)^2+1=41/16$ (two and one) or $(3/2)^2$ (all
in one). Hadamard's inequality gives
\begin{equation}\label{eq:hadamard}
 0<|\det D_J|\le2^{\sum_{i\in J}b_i}3^{q/2} .
\end{equation}
If $q<c$, or $q=c$ and $f>0$, then by \eqref{eq:fullcomp},
\eqref{eq:minorfactor} and \eqref{eq:hadamard},
\[
 v_2(\det K_I)\le\sum_{\rm full}(v_i-1)+\sum_{\rm def}e_i+\sum_{i\in J}b_i
 +\tfrac q2\log_23
 =n-f-c-\!\!\sum_{\substack{{\rm def}\ i\notin J}}\!\!b_i+\tfrac q2\log_23<n-f-c+q,
\]
strictly since $q>0$ and $\sqrt3<2$; the right side is at most $n-1$, so
$v_2(\det K_I)\le n-2$. If $q=c$ and $f=0$, then $K$ is square and every
column has coordinate sum three, so $\one^tK=3\one^t$ and $3\mid\det K$, hence
$3\mid\det D$; with $\sum_i(e_i+b_i)=n-c$,
\[
 2^{v_2(\det K)}\le2^{\sum e_i}\frac{|\det D|}3\le2^{n-c}3^{c/2-1}<\frac{2^n}3
 <2^{n-1}.
\]
In both cases \eqref{eq:cramer} gives $\ord(\det A)\le2^{n-2}$. So equality
occurs only with $q=0$. Finally $F_n$ is fixed by
$\diag(\xi,\xi^{-2},\dots,\xi^{(-2)^{n-1}})$ with $\xi$ of order $2^{n-1}$,
whose determinant $\xi^{S}$, $S=\sum_{i<n}(-2)^i$, has order $2^{n-1}$ since
$S$ is odd.
\end{proof}

